\documentclass[10pt,a4paper]{amsart}
\usepackage[margin=19mm]{geometry}
\usepackage{amsmath,amssymb,mathtools}
\usepackage[scr=rsfs,cal=euler]{mathalfa}
\usepackage{xfrac}
\usepackage[unicode,hidelinks,bookmarks=false]{hyperref}
\newcommand{\ie}{i.e.\ }
\newcommand{\cf}{cf.\ }
\newcommand{\resp}{respectively\ }
\newcommand{\Subsection}{\subsection}
\providecommand{\nobreakdash}{\nobreak\hbox{-}}
\setlength{\emergencystretch}{2em}
\multlinegap0pt
\usepackage{amssymb}
\usepackage{mathtools}
\usepackage{dsfont}
\usepackage{bbm}
\usepackage{subfig}
\usepackage{enumerate}
\usepackage{xcolor}
\newtheorem{lemma}{Lemma}
\newcommand{\N}{\mathbb{N}}
\newcommand{\di}{\mathrm{d}}   % differential
\title{Linearised Calder\'on problem: Reconstruction of unbounded perturbations in 3D}
\author{Henrik Garde, Markus Hirvensalo}
\date{Source: arXiv:2403.16588v1 (CC BY 4.0)}
\begin{document}
\maketitle

\noindent\emph{Source and licence.} Linearised Calder\'on problem: Reconstruction of unbounded perturbations in 3D. Version arXiv:2403.16588v1 (CC BY 4.0), distributed by arXiv under the Creative Commons Attribution 4.0 International licence, http://creativecommons.org/licenses/by/4.0/, as recorded in the arXiv licence metadata for this exact version and shown on the abstract page https://arxiv.org/abs/2403.16588v1. Full licence notice: \texttt{licenses/arxiv-2403.16588-CC-BY-4.0.txt}.\par\smallskip

\noindent\textit{Editorial scope.} Counted unit: the article's lemma lemma:3gradient (source0.tex lines 371--378). The article's complete original proof of it is delivered with it (source0.tex lines 380--427, one \texttt{proof} environment of 48 lines). No mathematics is written, repaired or re-derived: the statement and the proof are the article's own lines, in the article's own notation, and no retained line is substituted or re-flowed.  Retained uncounted as the source-local context the proof needs: lines 325--327.  Nothing was omitted from any retained span, so the package carries no omission record.  No retained line contains a citation, so the wrapper carries no bibliography and no bibliography entry.  No retained line contains a figure environment, an image-inclusion command or drawing code, so no image is delivered and the package declares no asset dependency.  Editorial formatting only, in addition: an AMS \texttt{amsart} preamble that restates the article's own declarations and macros used by the retained lines, namely the environments \texttt{lemma} on separate counters, together with the standard packages those lines need.  The retained environments print in this document as Lemma 1 in order of appearance, whereas the article prints the same items with the numbering of its own full document.  The complete native source of the article as distributed in the arXiv e-print for this exact version is bundled unchanged as \texttt{sources/156588/source0.tex}.
\begin{equation} \label{eq:LBeig}
    \Delta_{\partial B}g = -\ell(\ell+1)g, \quad g\in\mathcal{H}_\ell.
\end{equation}

\begin{lemma} \label{lemma:3gradient}
    Let $\ell_j\in \N_0$ and $g_{\ell_j} \in \mathcal{H}_{\ell_j}$ for $j\in\{1,2,3\}$. Then
    %
    \begin{equation*}
        \int_{\partial B} (\nabla_{\partial B} g_{\ell_1} \cdot \nabla_{\partial B} g_{\ell_2}) g_{\ell_3} \, \di S = \frac{\ell_1 (\ell_1 + 1) + \ell_2 (\ell_2 + 1) - \ell_3 (\ell_3 + 1)}{2} \int_{\partial B} g_{\ell_1} g_{\ell_2} g_{\ell_3} \, \di S.
    \end{equation*}
    %
\end{lemma}

\begin{proof}
    Below $\nabla$ denotes a gradient in $B$. We extend $g_{\ell_1}$, $g_{\ell_2}$, and $g_{\ell_3}$ to $B\setminus\{0\}$, constant in the radial direction. Using spherical coordinates gives
    %
    \begin{equation} \label{eq:gip}
        \nabla g_{\ell_i}\cdot\nabla g_{\ell_j} = r^{-2}\nabla_{\partial B} g_{\ell_i}\cdot \nabla_{\partial B} g_{\ell_j},
    \end{equation}
    %
    and, using \eqref{eq:LBeig},
    %
    \begin{equation} \label{eq:deltag}
        \Delta g_{\ell_j} = r^{-2} \Delta_{\partial B} g_{\ell_j} = -r^{-2}\ell_j (\ell_j + 1) g_{\ell_j}.
    \end{equation}
    %
    In the computations below, one can replace $B$ by $B\setminus B_\epsilon$ for an origin-centered ball $B_\epsilon$ with radius $\epsilon>0$, and let $\epsilon\to 0$. Since the considered functions are integrable on $B$, the limiting process gives precisely the results below, so we avoid these unnecessary technicalities in what follows. 
    	
    Define
    %
    \begin{equation} \label{eq:Edef}
        E_{\ell_1,\ell_2,\ell_3} = \int_{B} (\nabla g_{\ell_1} \cdot \nabla g_{\ell_2}) g_{\ell_3} \, \di x.
    \end{equation}
    %
    Consider the first integral on the right hand-side below: Integrating by parts, noting that the boundary term vanishes due to a vanishing normal derivative of $g_{\ell_2}$, and using the product rule, gives
    %
    \begin{equation} \label{eq:Ecyclic}
        E_{\ell_1,\ell_2,\ell_3} = -\int_B g_{\ell_1}(\Delta g_{\ell_2})g_{\ell_3}\,\di x - E_{\ell_2,\ell_3,\ell_1}.
    \end{equation}
    %
    Note that the assumptions needed for arriving at \eqref{eq:Ecyclic} are symmetric in terms of the indices $\ell_1,\ell_2,\ell_3$, and they can therefore be permuted to arrive at other such formulas. Thus we may use the cyclic property \eqref{eq:Ecyclic} in the following way:
    %
    \begin{align}
        E_{\ell_1,\ell_2,\ell_3} &= -\int_B g_{\ell_1}(\Delta g_{\ell_2})g_{\ell_3}\,\di x - E_{\ell_2,\ell_3,\ell_1} \notag \\
   		&= -\int_B g_{\ell_1}(\Delta g_{\ell_2})g_{\ell_3}\,\di x + \int_B g_{\ell_2}(\Delta g_{\ell_3})g_{\ell_1}\,\di x + E_{\ell_3,\ell_1,\ell_2} \notag \\
   		&= -\int_B g_{\ell_1}(\Delta g_{\ell_2})g_{\ell_3}\,\di x + \int_B g_{\ell_2}(\Delta g_{\ell_3})g_{\ell_1}\,\di x - \int_B g_{\ell_3}(\Delta g_{\ell_1})g_{\ell_2}\,\di x - E_{\ell_1,\ell_2,\ell_3}. \label{eq:Eiter}
    \end{align}
    %
    Combining \eqref{eq:deltag} and \eqref{eq:Eiter} gives
    %
    \begin{align*}
        E_{\ell_1,\ell_2,\ell_3} &= \frac{\ell_1 (\ell_1 + 1) + \ell_2 (\ell_2 + 1) - \ell_3 (\ell_3 + 1)}{2} \int_{B} r^{-2}g_{\ell_1} g_{\ell_2} g_{\ell_3} \, \di x \\
   		&= \frac{\ell_1 (\ell_1 + 1) + \ell_2 (\ell_2 + 1) - \ell_3 (\ell_3 + 1)}{2} \int_{\partial B} g_{\ell_1} g_{\ell_2} g_{\ell_3} \, \di S.
    \end{align*}    	
    %
    The proof is concluded by using \eqref{eq:gip} and \eqref{eq:Edef}:
    %
    \begin{equation*}
        E_{\ell_1,\ell_2,\ell_3} = \int_{\partial B} (\nabla_{\partial B} g_{\ell_1} \cdot \nabla_{\partial B} g_{\ell_2}) g_{\ell_3} \, \di S. \qedhere
    \end{equation*}
\end{proof}

\end{document}
